\chapter[Causal Chain Derivation Method]{Method for Causal Chain Derivation of Safety Performance Indicators}
\label{ch:methodology}

This chapter presents the core contribution of the thesis, a systematic Causal Chain
Derivation Method (CCDM) that produces calibrated leading safety performance indicators
from lagging ones and attaches them to a \ac{GSN} safety case. As established in
Chapter~\ref{ch:sota}, the reasoning components already exist in separate literatures
but are not integrated into a reproducible procedure. The method combines three of
them, namely the top-down decomposition of Fault Tree Analysis \cite{iec61025}, the
counterfactual prevention test of Why-Because Analysis \cite{ladkin1998}, and
threshold calibration grounded in precursor-estimation theory \cite{bieryi1995}. The
derived indicators are modelled in a \ac{GSN} modelling tool and evaluated on
simulation data generated in the CARLA simulator \cite{dosovitskiy2017carla}. The chapter first gives
an overview of the seven steps, then describes each step, presents worked examples,
and closes with the properties of the method that carry over into its instantiation
in Chapter~\ref{ch:experiment}.

\section{Overview of the Causal Chain Derivation Method}
\label{sec:method_overview}

The method is defined to be use-case-agnostic, so that it applies without change to a
left-turn, construction-zone, or pedestrian-crossing use case, or to any future one.
It consists of a preparatory step and seven derivation steps, of which
Steps~2 to~7 form a loop executed once for every lagging indicator in the register
built in Step~1. The complete sequence is shown in
Figure~\ref{fig:derivation_flow}. The method is architecture-independent in its
reasoning, because the subsystem decomposition in Step~3 names the functions that any
driving system must perform rather than the components that a particular
implementation contains. Architecture dependence enters only at Step~4, where a
required signal may or may not be exposed by a given implementation, and this
distinction gives rise to the two indicator tiers described in
Section~\ref{sec:method_signal}.

% ---------------------------------------------------------------------------
\begin{figure}[htbp]
    \centering
    \includegraphics[width=13cm]{spi_derivation_flow}
    \caption{The seven-step causal-chain derivation method. Steps~2 to~7 are executed
    once for each lagging safety performance indicator in the register produced by
    Step~1, and the loop terminates when the register is exhausted.}
    \label{fig:derivation_flow}
\end{figure}
% ---------------------------------------------------------------------------

\section{Use Case and Lagging-Indicator Register}
\label{sec:method_register}

The preparatory step, Step~0, defines the use case in terms of the function the system
performs, the actors present, and the operational design domain. Step~1 then builds
the lagging-indicator register by listing every loss event for the use case as a
lagging \ac{SPI}, each with a single loss-event type and a single unit of measurement.
This register is the set over which the derivation loop iterates. Candidate loss events
are drawn from a catalogue of twelve families, ranging from collision and contact events,
through near-miss and last-resort response events, to control-handover and operating-
domain violations. Near-miss and last-resort response events are the preferred
starting point, because they occur more frequently than collisions yet carry the same
causal information, which mirrors the precursor rationale of
Section~\ref{sec:sota_reactive_proactive}. A selection rule governs the register: each
lagging indicator must resolve to exactly one proximate cause in Step~2, and a
candidate that resolves to several unrelated causes is defined too broadly and is
split before proceeding.

\section{Proximate Cause by the Prevention Test}
\label{sec:method_proximate}

Step~2 identifies the proximate cause of each lagging indicator by applying the
prevention test, which asks whether the loss event would have been prevented if one
condition had been corrected at the last possible moment. This test is the
counterfactual test of Why-Because Analysis applied prospectively
\cite{ladkin1998}. The answer is expressed as a functional impossibility in passive
form, of the type that the system could not perform a required action, and exactly one
proximate cause is admitted per lagging indicator. The boundary between a proximate
and an intermediate cause is set by whether a response window remained after the
failure. If a detection failure left no time for any response, the detection failure
is the proximate cause, whereas if a response window remained and was not used, the
detection failure is an intermediate cause and the resulting avoidance failure is the
proximate cause.

\section{Causal Decomposition with Gate Logic}
\label{sec:method_decomposition}

Step~3 decomposes the proximate cause into intermediate causes using the gate logic of
Fault Tree Analysis \cite{iec61025}. An OR gate applies when any one cause is
sufficient on its own, which requires a separate leading indicator for each branch and
is the common case for perception and planning failures. An AND gate applies when the
causes must co-occur because redundancy is present, in which case a single shared
indicator can suffice. The gate for a pair of candidate causes is determined by asking
whether either cause alone can produce the proximate cause without the other. The
decomposition terminates when each cause maps to exactly one functional subsystem,
namely perception, prediction, planning, control, or operating-domain monitoring, and
is describable by a quantity that any implementation of that subsystem would produce.
Deeper decomposition is deliberately avoided, because it reaches algorithm-specific or
sensor-specific detail and breaks the implementation-agnostic character of the method.
A supplementary System-Theoretic Process Analysis cross-check is recommended for
control-loop and non-failure hazards that Fault Tree Analysis does not capture under
the safety of the intended functionality \cite{iso21448,stpahandbook}.

The use of gate logic requires reconciliation with the criticism of Fault Tree Analysis
made in Chapter~\ref{ch:sota}, namely that binary component-failure logic does not
represent the continuous, non-failure insufficiencies characteristic of a
machine-learning perception stack. The reconciliation rests on what the gates are
applied to. The decomposition does not enumerate component failure states; it partitions
a proximate cause into functional causes, each of which states that a required function
was not delivered adequately, and adequacy is fixed by the threshold calibrated in
Step~6 rather than by a binary health state. A gate therefore records which functional
shortfalls are jointly or severally sufficient to produce the proximate cause, while the
continuous quantity is carried by the signal and its threshold rather than by the gate.
This preserves the deductive structure that makes the derivation reproducible and places
the continuity where it belongs. One limitation remains and is stated rather than
concealed: a cause arising only from the interaction of two functions, neither of which
is individually inadequate, is not expressible in this logic, and the System-Theoretic
Process Analysis cross-check is the intended remedy for that class of cause.

\section{Observable Signal and Observability Check}
\label{sec:method_signal}

Step~4 names, for each intermediate cause, the earliest measurable quantity with units
that tracks the development of that cause before it fully manifests, together with the
lead time by which it does so. Step~5 then subjects each candidate signal to two tests
that both must pass. The first is computability, meaning that the signal can be
computed at runtime or from available logged data at a sufficient frequency. The
second is independence, meaning that the source of the signal is not the subsystem
being monitored, since a degraded subsystem that reports on itself would appear
healthy at exactly the moment monitoring is required. A signal that fails either test
is unusable regardless of how well it correlates with the loss event, and the requirement to
state signal provenance explicitly is claimed as a contribution, because published
indicator sets rarely record it.

The independence claimed for the behavioural tier requires one qualification, because it
governs how the results of this work transfer to a deployed system. In the simulation
environment of Section~\ref{sec:method_carla} the kinematic quantities are read from
ground truth, which is external to the vehicle software, so a behavioural-tier indicator
is independent of the driving function in the strict sense. In a deployed system no
ground truth exists, and the same quantities are estimated on board by perception and
localisation. A degradation of that estimation chain corrupts the indicator and the
driving function together, so the independence that holds by construction in simulation
becomes an assumption that must be argued separately in deployment, for example by
deriving the indicator from a sensing path not shared with the function it monitors. The
behavioural tier is independent of the planning and control functions it monitors in
both settings, but independent of the perception function only in simulation, and the
results reported for that tier are to be read under this restriction.

The independence requirement partitions the derivable indicators into two tiers,
summarised in Table~\ref{tab:tiers}. Behavioural-tier indicators are computed purely
from vehicle kinematics and scene state and are portable across any architecture,
including an opaque end-to-end network or a human driver. Architectural-tier indicators
require a signal exposed from inside a subsystem and therefore presuppose a modular
pipeline. The trade-off is explicit: behavioural-tier indicators buy portability but
give lower diagnostic resolution, since they reveal that a margin is eroding rather
than which subsystem caused it, and they typically provide shorter lead time, since
internal degradation usually precedes its kinematic consequence.

% ---------------------------------------------------------------------------
\begin{table}[htbp]
    \centering
    \caption{The two indicator tiers produced by the observability and independence
    check of Step~5.}
    \label{tab:tiers}
    \begin{tabular}{l l l}
        \toprule
        \textbf{Property} & \textbf{Behavioural tier} & \textbf{Architectural tier} \\
        \midrule
        Signal source     & vehicle kinematics, scene state & internal subsystem output \\
        Examples          & time-to-collision, closing speed, & detection confidence, \\
                          & required deceleration, PET & time-of-first-detection \\
        Architecture      & portable across any design & modular pipeline required \\
        Diagnostic resolution & margin erosion only & subsystem attribution \\
        Typical lead time & shorter & longer \\
        \bottomrule
    \end{tabular}
\end{table}
% ---------------------------------------------------------------------------

% =============================================================================
% DROP-IN REPLACEMENT for Section 4.6 of Chapter 4 (Method).
% Paste this in place of the existing Section 4.6. It carries the label
% sec:method_thresholds, which Chapters 5, 6 and 7 reference.
% Addresses the fix list: ALARP, GAMAB and MEM named as not applied; T defined
% as the exact one-sided binomial limit at alpha = 0.01 from Campaigns 2 and 3;
% the percentile example corrected to the 10th and 95th.
% =============================================================================

\section{The Metric Threshold and the SPI Threshold}
\label{sec:method_thresholds}

Steps~1 to~5 produce a signal and a source. A signal alone is not yet an
indicator: a quantity becomes a safety performance indicator only once a
threshold decides when its value is to be treated as a violation, and a second
threshold decides when the frequency of such violations is to be treated as a
departure from the argued behaviour. Step~6 sets both, and the general form of an
SPI used throughout this work is given in Equation~(\ref{eq:method_spiform}).

\begin{equation}
    \left\{
        \begin{array}{c}
            \text{No. of times the metric crosses } \theta
            \text{ in the unsafe direction} \\
            \text{per 100k manoeuvres, km or h driven}
        \end{array}
    \right\} < T
    \label{eq:method_spiform}
\end{equation}
\equationentry{General form of a safety performance indicator}

\textbf{Variable key:} $\theta$, metric threshold, in the unit of the metric; T,
SPI threshold, as a count per exposure unit.

The two thresholds answer different questions and are set by different means.
The metric threshold decides a single run and is calibrated from observed
behaviour; the SPI threshold bounds the violation count of a population of runs
and is set statistically. Neither is a risk-acceptance criterion.

\subsection{Setting the Metric Threshold}
\label{sec:method_metricthreshold}

Five ways of setting a metric threshold were considered before the method fixed
on one, and the choice is argued here rather than assumed, because it determines
whether the later evaluation is circular.
Table~\ref{tab:threshold_methods} compares the five by the data and the
assumptions each requires. The comparison is qualitative: no numerical benchmark
between the five was run, and none is implied.

% ---------------------------------------------------------------------------
% Copied verbatim from tables.tex
\begin{table}[htbp]
\centering
\caption{Options for the metric threshold. Five ways to set the metric threshold $\theta$ compared by the data and assumptions they require.}
\label{tab:threshold_methods}
\footnotesize
\setlength{\tabcolsep}{4pt}
\begin{tabular}{@{}>{\raggedright\arraybackslash}p{2.8cm}>{\raggedright\arraybackslash}p{2.1cm}>{\raggedright\arraybackslash}p{3.8cm}>{\raggedright\arraybackslash}p{3.4cm}>{\raggedright\arraybackslash}p{2.3cm}@{}}
\toprule
{\bfseries Method} & {\bfseries Needs collision labels} & {\bfseries Main assumption} & {\bfseries Data required} & {\bfseries Decision} \\
\midrule
Empirical percentile & no & uneventful runs represent normal driving & collision-free runs (available) & used \\
ROC curve, Youden index & yes & labelled collisions are representative & labelled collisions (the evaluation set) & not used: circular \\
Extreme value theory & no & a parametric tail model holds & long tails of the metric & future work \\
Kinematic bounds & no & friction and actuator limits are known per run & road and vehicle parameters & not used \\
Human-driver baselines & no & careful human driving is the reference & naturalistic data (inD, SIND) & future work \\
\bottomrule
\end{tabular}
\par\smallskip\footnotesize\raggedright Qualitative comparison by requirements, not a numerical benchmark.
\end{table}

The \textbf{empirical percentile} sets $\theta$ to a fixed percentile of the
per-run metric over runs that show normal behaviour. It requires no collision
labels, no distributional assumption and no vehicle parameters, and it is
reproducible from the data alone. Its cost is explicit: a tenth-percentile
threshold is crossed by about one tenth of uneventful runs by construction, so
the threshold marks unusual behaviour rather than unacceptable behaviour.

\textbf{ROC analysis with Youden's index} selects the threshold that maximises
sensitivity plus specificity on labelled data and so gives the best available
trade-off between false alarms and missed collisions \cite{youden1950}. It requires collision labels. The only labels available here are the
collisions against which the indicators are evaluated, so the threshold would be
fitted on the evaluation set and the evaluation would become circular.

\textbf{Extreme value theory} fits a parametric model to the tail of the metric
distribution and can therefore link $\theta$ to a target collision probability
\cite{songchitruksa2006evt}. It requires a
tail-model assumption and enough extreme observations to fit it, which is a study
in its own right.

\textbf{Kinematic bounds} derive $\theta$ from physical limits, for instance from
the deceleration the vehicle can actually achieve. Such bounds are explainable,
but they require friction, actuator and road-state assumptions for every run.
Those quantities vary with the sampled weather of
Section~\ref{sec:exp_parameters} and are not logged as ground truth.

\textbf{Human-driver baselines} set $\theta$ to match the behaviour of careful
human drivers, as an industry best practice describes \cite{avsc2023human}. This requires
naturalistic driving data such as the inD and SinD trajectory datasets
\cite{bock2020ind} \cite{xu2022sind}, which were not
evaluated in this work.

Three criteria decide between them. First, independence from the evaluation: the
threshold must not be fitted on the collisions against which it is tested, which
excludes the ROC criterion. Second, assumptions: the percentile needs none beyond
the choice of the calibration population, whereas extreme value theory needs a
tail model and kinematic bounds need per-run physical parameters. Third,
availability of data within this work: collision-free simulation runs exist in
large numbers, whereas naturalistic data and per-run physical parameters do not.
The empirical percentile is the only option that satisfies all three, and it is
the one used.

The threshold is computed from one value per run, namely the minimum or the
maximum of the metric over the run, because consecutive simulation steps are
strongly autocorrelated and would otherwise be counted as independent
observations. The calibration population is the uneventful runs of Campaign~1,
that is the runs without a collision and without an emergency braking event. The
tenth percentile is taken for the metrics in which small values are dangerous and
the ninety-fifth for REQ\_DECEL, in which large values are dangerous. Bootstrap
intervals confirm that the resulting values are stable;
Section~\ref{sec:exp_calibration} reports them. The width of the two tails is
itself a choice, and a sensitivity sweep over the percentile remains open.

The threshold depends on the configuration from which it is calibrated. Between
Campaign~1 and the realistic-driving campaigns the threshold of REQ\_DECEL shifts
by 21.5\,\% and that of PED\_CPA by 37.9\,\%. It is therefore calibrated once and
applied unchanged to every campaign rather than refitted per campaign, so that
the same decision rule is applied to every run of the evaluation.

\subsection{Setting the SPI Threshold}
\label{sec:method_spithreshold}

The metric threshold leaves open how many violations are too many. The SPI
threshold answers that, and it is set statistically because the inputs for a
risk-based answer are not available in this work.

T is the exact one-sided binomial limit at a significance level of
$\alpha = 0.01$ on the violation rate observed in Campaigns~2 and~3, evaluated
for the number of eligible runs in the window being checked. It is the smallest
violation count whose probability under that baseline rate falls to or below
$\alpha$; Equation~(\ref{eq:exp_spithreshold}) in
Section~\ref{sec:exp_criteria} states it formally. The rate is checked over
monitoring windows of 200 runs. The window is the unit over which the check is
performed and not the unit of T, which is stated per exposure unit in
Equation~(\ref{eq:method_spiform}).

A threshold set this way bounds departure from a measured baseline. Reaching T
means that the violation count of a window has become improbable under the
behaviour the baseline campaigns demonstrated; it does not mean that a level of
risk has been exceeded, because no level of risk has been stated. A
risk-acceptance SPI threshold would instead be derived from a principle such as
ALARP \cite{hse2001r2p2}, GAMAB, which EN~50126-2 terms GAME, or MEM
\cite{en50126}, or from the benchmark of a careful and competent human driver
\cite{wang2024sotif} \cite{avsc2023human}. Those principles are named
here to place the statistical threshold against them; none of them is applied in
this work, and setting T from one of them is left to
Section~\ref{sec:outlook}.

\section{Attachment to the Safety Case}
\label{sec:method_gsn}

Step~7 attaches each derived leading indicator to the \ac{GSN} goal that it monitors,
attaches the lagging parent indicator to the same goal, records $\theta$ and $T$ as
dimension elements, and documents the full derivation chain in the traceability field
of the model in the \ac{GSN} modelling tool. This closes the loop opened in Step~1, since the goal that
began as a bare assertion now carries a monitorable claim, and a violation invalidates
that specific claim rather than leaving the safety case silently wrong. Two failure
modes are avoided explicitly, namely attaching an indicator to a convenient parent
goal instead of the goal it actually monitors, and omitting the traceability record,
which is the element that distinguishes a derived indicator from an asserted one. The
resulting argument structure is consistent with the \ac{GSN} community standard
\cite{gsncs2021} and with the argument pattern for the use of indicators proposed by
Ratiu et al. \cite{ratiu2024}. The method does not depend on any particular tool; a
model-based safety-engineering environment that supports \ac{GSN}, such as SafeTbox, is
recommended for carrying out the attachment in practice.

\section{Worked Examples}
\label{sec:method_examples}

The method is illustrated on three loss events drawn from the target use cases, one per
use case, and summarised in Table~\ref{tab:examples}. The three cases are given to
show that the derivation is use-case-agnostic, which is a claim of the method rather than
of the evaluation; only the left-turn case is carried into the empirical work of
Chapter~\ref{ch:results}, in line with the scope set in
Section~\ref{sec:intro_scope}. Each example begins from a lagging
indicator, identifies a single proximate cause, decomposes it with an OR gate into two
intermediate causes attributed to distinct subsystems, and yields one leading indicator
per branch. The left-turn example is developed in full in
Figure~\ref{fig:causal_chain_example}. Its lagging indicator is the rate of collisions
with an oncoming vehicle, its proximate cause is that the ego vehicle occupied the
conflict area while an oncoming vehicle was closing on it and no avoidance action
succeeded in time, and its two branches are a gap that was already insufficient at the
moment of commitment, attributed to planning, and a closing speed that exceeded what
the turn arc allowed, attributed to prediction. The example also illustrates the
placement of the reaction tier. An emergency-braking activation with an oncoming
vehicle inside the collision horizon is not a lagging indicator, because no loss has
yet occurred; it is an intermediate leading indicator, the last-resort mitigation
nearest the loss, and it is therefore positioned above the kinematic decomposition and
attached to the safety case independently of it rather than being derived from the
kinematic branches. The two kinematic branches
carry different lead times, of the order of three to eight seconds and of half a second
to three seconds respectively, which demonstrates the spanning property described in
Section~\ref{sec:method_thresholds}.

% ---------------------------------------------------------------------------
\begin{table}[htbp]
    \centering
    \caption{Three worked derivations, one per use case, from a lagging indicator
    through its proximate cause to a pair of leading indicators, each attributed to the
    functional subsystem in which its cause originates.}
    \label{tab:examples}
    \small
    \begin{tabular}{@{}p{2.6cm} p{5.4cm} p{5.6cm}@{}}
        \toprule
        \textbf{Use case} & \textbf{Lagging indicator and proximate cause} & \textbf{Leading indicators (subsystem)} \\
        \midrule
        Left turn & collision with an oncoming vehicle; the
        conflict area was occupied while an oncoming vehicle was closing and no
        avoidance succeeded in time & gap TTC at commitment (planning);
        required deceleration over the arc (prediction) \\[2pt]
        Pedestrian crossing & collision with a crossing pedestrian; the turn was
        committed without time to stop & time to closest point of approach
        (planning); crosswalk occupancy confidence (perception) \\[2pt]
        Construction zone & collision with or near-miss of a worker; the system
        could not stop in time & TTC at first detection (perception); detection
        confidence per frame (perception) \\
        \bottomrule
    \end{tabular}
\end{table}
% ---------------------------------------------------------------------------

% ---------------------------------------------------------------------------
\begin{figure}[htbp]
    \centering
    \includegraphics[width=15cm]{spi_causal_chain_example}
    \caption{Worked causal-chain derivation for the left-turn use case. The lagging
    indicator, a collision with an oncoming vehicle, is decomposed through its proximate
    cause and an OR gate into two subsystem-level intermediate causes, each yielding a
    leading indicator with its own earliest observable signal and lead time. The
    emergency-braking activation is the intermediate leading indicator of the reaction
    tier and attaches to the safety case independently of the kinematic branches.}
    \label{fig:causal_chain_example}
\end{figure}
% ---------------------------------------------------------------------------
% PASTE INTO: Chapter 4, Section 4.8 (Worked Examples)
% WHERE:      directly after the existing Figure 4.2 (the ONC chain figure)
% UPLOAD:     spi_causal_chain_ped.pdf and spi_causal_chain_rear.pdf into your Graphics folder

% First add this sentence to the end of the paragraph above Figure 4.2:
%
%   The same derivation is carried out for the other two conflict chains of the
%   manoeuvre. Figure~\ref{fig:chain_ped} gives the pedestrian chain and
%   Figure~\ref{fig:chain_rear} the following chain. The three differ in where the
%   derivation terminates: the oncoming and pedestrian chains each retain a
%   reaction-tier indicator, whereas the following chain has none, and the
%   pedestrian and following chains each lose one branch at the observability and
%   independence check of Step~5.

\begin{figure}[htbp]
  \centering
  \includegraphics[width=15cm]{spi_causal_chain_ped}
  \caption{Worked causal-chain derivation for the pedestrian-crossing
    conflict. The lagging indicator, a collision with a crossing pedestrian,
    is decomposed through its proximate cause and an OR gate into two
    subsystem-level intermediate causes. Only the planning branch yields an
    indicator that survives Step~5; the perception branch is derived but is
    of the architectural tier and therefore outside the empirical scope of
    this work.}
  \label{fig:chain_ped}
\end{figure}

\begin{figure}[htbp]
  \centering
  \includegraphics[width=15cm]{spi_causal_chain_rear}
  \caption{Worked causal-chain derivation for the following-vehicle conflict.
    The chain carries no reaction-tier indicator, because the ego vehicle has
    no last-resort mitigation directed rearward, and the branch attributing
    the conflict to the follower's reaction delay is rejected at Step~5 as
    not observable to the ego vehicle.}
  \label{fig:chain_rear}
\end{figure}
\section{Instantiation of the Method}
\label{sec:method_carla}

The seven steps above are defined without reference to a particular use case or
simulator, which is what allows the same procedure to be applied to the three conflict
chains of a single manoeuvre. Their instantiation is the subject of
Chapter~\ref{ch:experiment}, which fixes the simulation environment and the ego
configuration, describes each of the three conflict chains and the mechanism by which
conflicts in it are generated, states the indicator register that Steps~3 to~7 produce
for those chains, and defines the criteria against which the resulting indicators are
judged. Chapter~\ref{ch:results} then reports the evaluation against those definitions.

Two properties of the instantiation are noted here because they follow from the method
rather than from the experiment. First, because the configuration used exposes no
perception internals, every indicator that survives the observability and independence
check of Section~\ref{sec:method_signal} is of the behavioural tier; the
architectural-tier indicators that depend on detection confidence or time of first
detection are outside the empirical scope of this work, since ground truth carries no
notion of confidence and a calibrated perception model would be required before any
threshold on such a signal is meaningful. Second, indicators are computed in
post-processing rather than during a run, which permits $\theta$ and $T$ to be
recalibrated without regenerating the scenarios, and which equally means that the
runtime cost of evaluating an indicator on a vehicle is not measured here. The
operational loop of Chapter~\ref{ch:sota} is accordingly evidenced in its derivation and
monitoring stages only.